\documentclass[10pt]{article}
\usepackage{compresr}
\cprSetHeader{TypeScript · LangGraph Integration}

\begin{document}

\cprTitleBlock{Compresr SDK · TypeScript}{LangGraph Integration}{%
    Compress graph state, checkpoints, and stores in LangGraph.js with
    drop-in, query-aware primitives. Cut token spend and storage size
    without rewriting your graph.}

% ===========================================================================
\section{Install}

\begin{shcode}
npm install @compresr/sdk
export COMPRESR_API_KEY=cmp_...
\end{shcode}

Peers (\cprInlineCode{@langchain/langgraph},
\cprInlineCode{@langchain/langgraph-checkpoint},
\cprInlineCode{@langchain/core}) are optional --- install only what you use.
Node $\geq$ 20, dual ESM+CJS, typed end-to-end. Import from the dedicated
entrypoint:

\begin{tscode}
import {
  makeCompresrNode, CompresrCheckpointSerializer, CompresrStore,
} from '@compresr/sdk/integrations/langgraph';
\end{tscode}

A \cprInlineCode{compresrHandoffTool} is also available for supervisor /
multi-agent patterns --- see the SDK reference.

% ===========================================================================
\section{Compression node}

\cprInlineCode{makeCompresrNode} returns an async graph node that reads a
state field, compresses it against a query, and returns a patch. No-ops when
input is missing, too short, or unchanged.

\begin{tscode}
const compress = makeCompresrNode<typeof State.State>({
  apiKey:                 process.env.COMPRESR_API_KEY!,
  contextKey:             'retrievedText',
  queryKey:               'question',
  targetCompressionRatio: 0.5,
  minTokens:              200,
});
graph.addNode('compress', compress);
\end{tscode}

% ===========================================================================
\section{Compressed checkpoints}

\cprInlineCode{CompresrCheckpointSerializer} plugs into the
\cprInlineCode{serde} slot of any LangGraph.js checkpointer
(\cprInlineCode{MemorySaver}, \cprInlineCode{SqliteSaver},
\cprInlineCode{PostgresSaver}). Long strings on the listed fields are
compressed on write. \textbf{Compression is lossy and one-way} — reads
return the compressed string, not the original. Use the
\cprInlineCode{fields} allow-list to scope compression to retrieved
text, transcripts, and other paraphrasable payloads; leave titles, IDs,
and anything you need to read back verbatim out of the list.

\begin{tscode}
const checkpointer = new MemorySaver(
  new CompresrCheckpointSerializer({
    apiKey:                 process.env.COMPRESR_API_KEY!,
    fields:                 new Set(['retrievedText', 'transcript']),
    minTokens:              500,
    targetCompressionRatio: 0.5,
  }),
);
const graph = builder.compile({ checkpointer });
\end{tscode}

Without \cprInlineCode{fields}, every string $\geq$
\cprInlineCode{minTokens} anywhere in the state tree is compressed
(including IDs and titles that exceed the threshold) — set the allow-list
in any production deployment.

% ===========================================================================
\section{Compressed store}

\cprInlineCode{CompresrStore} wraps any LangGraph.js
\cprInlineCode{BaseStore} by composition: reads and namespace operations
pass through; \cprInlineCode{put} and \cprInlineCode{batch} compress
allow-listed string fields above the token threshold.

\begin{tscode}
const store = new CompresrStore(new InMemoryStore(), {
  apiKey:    process.env.COMPRESR_API_KEY!,
  fields:    new Set(['retrievedText']),
  minTokens: 500,
});
await store.put(['users', 'u_42', 'memories'], 'memo_01', {
  retrievedText: '<3KB retrieval output>',
  score:         0.91,
});
\end{tscode}

% ===========================================================================
\section{Key parameters}

\renewcommand{\arraystretch}{1.15}
\begin{longtable}{@{}p{3.6cm} p{2.4cm} p{8.0cm}@{}}
\toprule
\textbf{Option} & \textbf{Default} & \textbf{Description} \\
\midrule
\endhead
\cprInlineCode{apiKey} / \cprInlineCode{client} & --- &
API key, or a pre-built \cprInlineCode{CompressionClient}. One required. \\
\cprInlineCode{contextKey} & --- &
\textit{(node only)} State key holding the text to compress. Required. \\
\cprInlineCode{fields} & --- &
\textit{(serializer / store)} Allow-list of parent keys to compress.
Strongly recommended. \\
\cprInlineCode{compressionModel} & \cprInlineCode{latte\_v1} &
\cprInlineCode{latte\_v2} is the faster variant. \\
\cprInlineCode{targetCompressionRatio} & \cprInlineCode{0.5} &
$[0,1]$ = fraction of tokens to remove; $>1$ = Nx factor. \\
\cprInlineCode{minTokens} & node: \cprInlineCode{200} \newline
serializer / store: \cprInlineCode{500} &
Skip compression below this threshold. Storage paths use a higher
default to avoid compressing short rows. \\
\cprInlineCode{coarse} & (server) &
\cprInlineCode{true} paragraph-level, \cprInlineCode{false} token-level. \\
\cprInlineCode{query} / \cprInlineCode{queryKey} / \cprInlineCode{queryExtractor}
  & --- &
Choose the query that drives compression: static string, state key, or
extractor function. Highest-priority non-empty wins. \\
\cprInlineCode{onError} & \cprInlineCode{passthrough} &
On transport, validation, or unexpected errors: return the original string
so the graph keeps running. Use \cprInlineCode{raise} to surface
exceptions. \\
\bottomrule
\end{longtable}

\end{document}
